\section{Lessons and Ongoing Deployment}
\label{sec:discussion}

\subsection{What the measurements establish}
\textbf{Reconstruction.} Restormer leads the six-model benchmark; RAMiT is
smaller but has no matched cascade/device measurements. All continuation
variants fail validation replacement, so epoch 51 remains selected.

\textbf{Detection and cost.} Local AP improves; Nano AP and throughput decline.
Different canvases, runtimes and scoring prevent attributing the gap to
resolution or precision alone. The five-object geometry diagnostic cannot
replace a full labelled resize-only Nano arm. Camera fidelity scores and
illustrative overlays do not establish object accuracy.

\subsection{Threats to validity}
\textbf{Data and training.} Prior test inspection, six smoke backgrounds,
imperfect scene independence, synthetic weather and differing pretraining limit
generalization. Single seeds leave training variability unmeasured; bootstrap
intervals condition on the checkpoints.

\textbf{Detection and device.} The pilot has 72 images; weather variants reuse
objects. Unlabelled camera counts are not AP. Short, single Nano trials do not
establish thermal stability. Live and replay timing boundaries differ, and
input-rail energy includes idle power.

\textbf{Compression.} Four fixtures cannot establish full-set task retention.
FP32-stored weight grids lack integer kernels; dense masks preserve tensor
shapes and file size. Rejection applies to the tested settings.

\subsection{Phase 7: RAMiT + YOLOv12 edge deployment}
\label{sec:ongoing}
As of 13 November 2026, this sole ongoing phase integrates RAMiT with YOLOv12.
Export checks precede matched labelled controls and repeated device measurements.
\textbf{No end-to-end accuracy or throughput is yet measured.} Broader validation
needs independent smoke scenes, multiple seeds, camera annotations and sustained
runs.

\subsection{Conclusion}
Restoration quality, detector accuracy and device cost require separate
measurements. The six completed phases establish these trade-offs and the
controls for the ongoing RAMiT + YOLOv12 deployment study.
